\section{Discussion}
\label{sec:discussion}

\paragraph{Behavioral signals as bias proxies.}
Our results provide evidence that user interaction patterns on social news platforms carry information about perceived media bias, even when users are not explicitly evaluating bias.
The consistency of the signal across multiple feature families (karma, sentiment, temporal, engagement) and across a 16-year span strengthens this finding.
However, the modest predictive performance (AUC $0.642$) indicates that interaction features alone are insufficient for reliable bias classification.
This suggests their value lies in complementing content-based approaches, not replacing them.

\paragraph{Karma entropy as a controversy marker.}
The finding that biased articles have higher karma entropy is particularly informative.
Unlike simple metrics such as average karma or score, entropy captures the \textit{diversity} of community reactions.
A high-entropy karma distribution means that the community is divided: some comments are strongly upvoted while others are downvoted, reflecting genuine disagreement rather than uniform sentiment.
This aligns with theoretical accounts of media bias as a stimulus for divergent interpretations \cite{vallone1985hostile}.

\paragraph{Faster reactions, shorter threads.}
The faster initial reaction times and shorter thread lifespans for biased articles suggest that bias triggers impulsive engagement.
Users respond quickly to biased content, but the conversation degrades faster (steeper karma drift) and dies sooner.
This pattern is consistent with the ``outrage engagement'' hypothesis \cite{rathje2021out}, where emotionally charged content drives rapid but shallow interaction.

\paragraph{Two media ecosystems.}
The detection of two distinct outlet communities, separated by commenter bases rather than editorial stance, reveals structural segregation in the Spanish online news ecosystem.
Community~1 (traditional outlets) and Community~2 (digital-native outlets) share relatively few active commenters, despite both communities containing outlets across the political spectrum.
This suggests that audience segregation on Men\'eame is driven more by platform habits and outlet format than by ideology alone.

\paragraph{User diversity and bias exposure.}
The finding that users with mixed bias exposure ($0.4$--$0.6$) read the most diverse set of outlets is notable.
It suggests that moderate bias exposure is not a sign of echo chamber membership, but rather of broad engagement with the news ecosystem.
Conversely, users at both extremes (low and very high exposure) tend to concentrate on fewer outlets.

\paragraph{Limitations.}
Several limitations should be noted.
First, our bias labels are automatic, derived from a classifier trained on a relatively small Spanish dataset (MBBMD, 100 articles).
We attempted cross-lingual transfer using a DistilBERT model trained on BEADS (English) and MBBMD, but the resulting model labeled only 2\% of articles as biased, indicating poor transfer.
The original Spanish-specific model, while imperfect, provides more plausible label distributions for this corpus.
Second, Men\'eame's user base is not representative of the general Spanish population: it skews male, tech-savvy, and politically engaged.
Findings may not generalize to other platforms or demographics.
Third, we lack comment threading information (reply-to relationships), which prevents analysis of discussion depth and argumentation structure.
Fourth, our study is correlational: we cannot determine whether bias causes different interaction patterns or whether confounding factors (topic, outlet popularity) drive both bias labels and engagement.
